% !TeX spellcheck = ru_RU
\documentclass[a4paper, 14pt, unknownkeysallowed]{extreport}
\usepackage[utf8]{inputenc}
\usepackage[T2A]{fontenc}
\usepackage[english,russian]{babel}
\usepackage{cmap}
\usepackage{enumitem}

\usepackage{csvsimple}
\usepackage{hyphenat} 
\usepackage{amstext, amsmath,amsfonts,amssymb,amsthm,mathtools} 

\usepackage{setspace}
\onehalfspacing 

\usepackage{geometry}
\geometry{left=30mm}
\geometry{right=10mm}
\geometry{top=10mm}
\geometry{bottom=20mm}

\usepackage{graphicx}
\graphicspath{{img/}} 

\usepackage{threeparttable}
\usepackage{bigdelim}

\usepackage{setspace}
\onehalfspacing % Полуторный интервал

\frenchspacing
\usepackage{indentfirst} % Красная строка

\usepackage{cases}

\usepackage[unicode,pdftex]{hyperref} % Ссылки в pdf
\hypersetup{hidelinks}

\usepackage{xcolor}
\usepackage{listings}
% Для листинга кода:
\lstset{%
	language=Matlab,   					% выбор языка для подсветки	
	basicstyle=\small\sffamily,			% размер и начертание шрифта для подсветки кода
	numbers=left,						% где поставить нумерацию строк (слева\справа)
	%numberstyle=,					% размер шрифта для номеров строк
	stepnumber=1,						% размер шага между двумя номерами строк
	numbersep=5pt,						% как далеко отстоят номера строк от подсвечиваемого кода
	frame=single,						% рисовать рамку вокруг кода
	tabsize=4,							% размер табуляции по умолчанию равен 4 пробелам
	captionpos=t,						% позиция заголовка вверху [t] или внизу [b]
	breaklines=true,					
	breakatwhitespace=true,				% переносить строки только если есть пробел
	escapeinside={\#*}{*)},				% если нужно добавить комментарии в коде
	backgroundcolor=\color{white},
}
\usepackage{caption}
\captionsetup{labelsep = endash}
\captionsetup[figure]{name = {Рисунок}, justification = centerlast}
\captionsetup[lstlisting]{justification = justified}
\captionsetup[table]{justification = justified}

\usepackage{titlesec}
\titleformat{\chapter}{\LARGE\bfseries}{\thechapter}{20pt}{\LARGE\bfseries}
\titleformat{\section}{\Large\bfseries}{\thesection}{20pt}{\Large\bfseries}

\renewcommand\labelitemi{\ --\ }
\renewcommand\labelenumi{\theenumi)}

\newcommand{\img}[3] {
	\begin{center}
		\begin{figure}[h!]
			\includegraphics[scale = #1]{img/#2}
			\caption{#3}
			\label{img:#2}
		\end{figure}
	\end{center}
}

\newcommand{\lst}[4]{
\lstinputlisting[language=Matlab, firstline=#1, lastline=#2, label=lst:#3,caption=#4]{../src/code.m}
}

\usepackage{siunitx,array,booktabs}
\usepackage{amsthm}
\usepackage{tabularx}
\newtheorem{theorem}{Theorem}
\newtheorem{definition}{Опредление}


\begin{document}
% !TeX spellcheck = ru_RU
\begin{titlepage}
	
	\newgeometry{pdftex, left=2cm, right=2cm, top=2.5cm, bottom=2.5cm}
	\fontsize{12pt}{12pt}\selectfont
	\noindent \begin{minipage}{0.15\textwidth}
		\includegraphics[width=\linewidth]{img/b_logo.jpg}
	\end{minipage}
	\noindent\begin{minipage}{0.9\textwidth}\centering
		\textbf{Министерство науки и высшего образования Российской Федерации}\\
		\textbf{Федеральное государственное бюджетное образовательное учреждение высшего образования}\\
		\textbf{«Московский государственный технический университет имени Н.~Э.~Баумана}\\
		\textbf{(национальный исследовательский университет)»}\\
		\textbf{(МГТУ им. Н.~Э.~Баумана)}
	\end{minipage}

	\noindent\rule{18cm}{3pt}
	\newline\newline
	\noindent ФАКУЛЬТЕТ $\underline{\text{«Информатика и системы управления»~~~~~~~~~~~~~~~~~~~~~~~~~~~~~~~~~~~~~~~~~~~~~~~~~~~~~~~}}$ \newline\newline
	\noindent КАФЕДРА $\underline{\text{«Программное обеспечение ЭВМ и информационные технологии»~~~~~~~~~~~~~~~~~~~~~~~}}$\newline\newline\newline\newline\newline\newline\newline
	
	
	\begin{center}
		\Large\textbf{Отчет по лабораторной работе № 4}
		
		\Large\textbf{по курсу <<Моделирование>>}
		
		\textbf{\newline}		
	\end{center}
	
	\noindent\textbf{Тема} $\underline{\text{~~Моделирование простейшей СМО~~~~~~~~~~~~~~~~~~~~~~~~~~~~~~~~~~~~~~~~~~~~~~~~~~~~~~~~~~~~~~~~~~~~~~~~}}$\newline\newline
	\noindent\textbf{Студент} $\underline{\text{~~Виноградов А. О.~~~~~~~~~~~~~~~~~~~~~~~~~~~~~~~~~~~~~~~~~~~~~~~~~~~~~~~~~~~~~~~~~~~~~~~~~~~~~~~~~~~~~~}}$\newline\newline
	\noindent\textbf{Группа} $\underline{\text{~~ИУ7-76Б~~~~~~~~~~~~~~~~~~~~~~~~~~~~~~~~~~~~~~~~~~~~~~~~~~~~~~~~~~~~~~~~~~~~~~~~~~~~~~~~~~~~~~~~~~~~~~~~~~~~}}$\newline\newline
	\noindent\textbf{Оценка (баллы)} $\underline{\textbf{~~~~~~~~~~~~~~~~~~~~~~~~~~~~~~~~~~~~~~~~~~~~~~~~~~~~~~~~~~~~~~~~~~~~~~~~~~~~~~~~~~~~~~~~}}$\newline\newline
	\noindent\textbf{Преподаватель} $\underline{\text{~~Рудаков И. В.~~~~~~~~~~~~~~~~~~~~~~~~~~~~~~~~~~~~~~~~~~~~~~~~~~~~~~~}}$\newline
	
	\begin{center}
		\vfill
		Москва~---~\the\year
		~г.
	\end{center}
	\restoregeometry
	
\end{titlepage}
\setcounter{page}{2}

\chapter{Теоретическая часть}
Для выполнения работы был выбран критерий «хи-квадрат». Это один из самых известных статистических критериев,  также это основной метод, используемый в сочетании с другими критериями. 
С помощью этого критерия можно узнать, удовлетворяет ли генератор случайных чисел требованию равномерного распределения или нет. 
Для оценки по этому критерию необходимо вычислить статистику V по формуле:

\begin{equation}
	\label{ass}
	V = \frac{1}{n}\sum_{s=1}^{k}(\frac{Y^{2}_S}{p_s}) - n,
\end{equation} 

где n – количество независимых испытаний, k – количество категорий, $Y_s$ — число наблюдений, которые действительно относятся к категории S,  $p_s$ — вероятность того, что случайное наблюдение относится к категории s.  

Значение V является значением критерия «хи-квадрат» для экспериментальных данных. Приемлемое значение этого критерия можно определить по таблице на рисунке \ref{img:1.jpg}. Для этого используем строку с v = k-1, где k = 10, 90, 900 для задания лабораторной.  P в этой таблице — это вероятность того, что экспериментальное значение $V_e$ будет меньше теоретического $V_t$ или равно ему. Ее также можно рассматривать как доверительную вероятность.

Если вычисленное по таблице $P_V$ окажется меньше 0.01 или больше 0.99, можно сделать вывод, что эти числа недостаточно случайные. Если $P_V$ лежит между 0.01 и 0.05 или между 0.95 и 0.99, то эти числа «подозрительны». Если $P_V$ лежит между 0.05 и 0.1 или 0.9-0.95, то числа можно считать «почти подозрительными». Обычно необходимо произвести проверку три раза и более с разными данными. Если по крайней мере два из трех результатов оказываются подозрительными, то числа рассматриваются как недостаточно случайные.  

\img{0.5}{1.jpg}{Некоторые процентные точки $\chi^2$- распределения.
	(Источник: Кнут Д. Э. «Искусство программирования» )
}

\chapter{Результат}
На рисунке \ref{img:ex} приведен пример работы программы.
\img{0.6}{ex}{Пример работы программы}


\end{document}